\chapter*{Sources}
\addcontentsline{toc}{chapter}{Sources}
\noindent Each chapter ends with the papers and books it uses. This list only points back.

\begin{description}
\item[Chapter~1] Mertens's theorems. Euler's E175, Mertens (1874), Chebyshev.
\item[Chapter~2] Averages of arithmetic functions. Wintner, Dirichlet's hyperbola, the divisor problem.
\item[Chapter~3] Characters and $L$-functions. Gauss sums, the non-vanishing of $L(1,\chi)$.
\item[Chapter~4] Dirichlet's theorem. The race modulo $4$, and the conditional bias of Rubinstein and Sarnak.
\item[Chapter~5] The prime number theorem, by Newman's contour. Apostol for the elementary equivalences; Titchmarsh, Chapter~14, for the quoted converse on $M(x)$.
\item[Chapter~6] The zeros of $\zeta$. Hadamard, von Mangoldt, de la Vallée Poussin. The proportion-of-zeros records are quoted in a remark.
\item[Chapter~7] The explicit formula. Von Mangoldt. The functional equation of $L(s,\chi)$ is admitted and proved in Volume~13, Chapter~12.
\item[Chapter~8] Sieves. Selberg, Brun, Brun--Titchmarsh. Montgomery and Vaughan's removal of the secondary term, and the sharp large-sieve constant, are quoted in remarks.
\item[Chapter~9] The Erd\H{o}s--Kac theorem. Kubilius's model.
\item[Chapter~10] Cramér's model. The Baker--Harman--Pintz exponent and the Ford--Green--Konyagin--Maynard--Tao gap are quoted in a remark. Maier's theorem is pointed at Volume~15.
\item[Chapter~11] Exponential sums. Van der Corput, Weyl, Graham and Ringrose (quoted).
\item[Chapter~12] Modular forms. Ramanujan, Mordell, Deligne (quoted). The closing problems have their own heading in the contents.
\end{description}
